\chapter{Atividades Futuras}

Como continuidade deste trabalho, as seguintes atividades estão previstas para o
TCC~II, abrangendo tanto a avaliação da ferramenta quanto a ampliação de suas
funcionalidades:

\begin{itemize}
  \item \textbf{Avaliação de desempenho} das importações com diferentes volumes de dados
        CSV e configurações, comparando importações simples e complexas para os diferentes
        modelos de dados destino.

  \item \textbf{Interface gráfica \textit{desktop}} (semestre seguinte) que monte o
        comando CLI e visualize os arquivos JSON de configuração.

  \item \textbf{Suporte a TSV/Excel} e a processamento \textit{chunked} para CSV maiores
        que a RAM.

  \item \textbf{Importação com um arquivo CSV por entidade.} Identificar a origem de
        cada dado é um requisito essencial do processo de importação: no formato atual,
        cada linha do CSV precisa indicar a entidade a que pertence, papel cumprido pela
        coluna discriminadora \texttt{action} em conjunto com os \texttt{filters} da
        configuração (Seção~\ref{sec:config-format}). Uma evolução prevista é aceitar um arquivo CSV por
        entidade, em que o próprio arquivo designa a origem dos dados, dispensando a
        coluna discriminadora e os filtros correspondentes; o CSV único combinado
        permaneceria como modo alternativo de entrada.

  \item \textbf{Operadores de filtro adicionais} e políticas de coerção de tipos
        configuráveis.

  \item \textbf{Testes de integração} contra \texttt{docker compose} em CI.

  \item \textbf{Generalização} do núcleo de importação para novos conectores
        (mensageria, \textit{data lakes}, outros SGBDs).
\end{itemize}
